\documentclass[14pt,a4paper,oneside]{extreport}
\usepackage{gost732}
\GostNumberWithinSections
\begin{document}
\begin{GostTitlePage}\begin{center}Проверка макета\end{center}\end{GostTitlePage}
\GostTOC
\chapter{Проверка первого раздела}
\section{Заголовок подраздела}
\subsection{Первый озаглавленный пункт}
Проверяются номера рисунка~\ref{fig:first}, формулы~\eqref{eq:first} и таблицы~\ref{tab:long}.
Дополнительные материалы приведены в приложениях~\ref{app:a} и~\ref{app:b}.
Текст со сноской\footnote{Пояснение на той же странице.}.
\begin{figure}[!htbp]\centering\fbox{Рисунок первого раздела}
\caption{Первый рисунок}\label{fig:first}\end{figure}
\begin{equation}\label{eq:first}a=b+c.\end{equation}
\GostWhere{$a$ — сумма;\\$b$ — первое слагаемое;\\$c$ — второе слагаемое.}
\subsection{Второй озаглавленный пункт}
\begin{GostDashes}\item первый элемент,\item второй элемент.\end{GostDashes}
\begin{GostNotes}\item Первое пояснение\item Второе пояснение\end{GostNotes}
\begin{longtable}{|p{75mm}|>{\centering\arraybackslash}p{75mm}|}
\caption{Результат наблюдения}\label{tab:long}\\
\hline Наблюдение & Значение\\\hline\endfirsthead
\multicolumn{2}{@{}l}{\GostSingle Продолжение таблицы \thetable}\\
\hline Наблюдение & Значение\\\hline\endhead
\hline\endfoot\hline\endlastfoot
Наблюдение 1 & 1\\
Наблюдение 2 & 2\\
Наблюдение 3 & 3\\
Наблюдение 4 & 4\\
Наблюдение 5 & 5\\
Наблюдение 6 & 6\\
Наблюдение 7 & 7\\
Наблюдение 8 & 8\\
Наблюдение 9 & 9\\
Наблюдение 10 & 10\\
Наблюдение 11 & 11\\
Наблюдение 12 & 12\\
Наблюдение 13 & 13\\
Наблюдение 14 & 14\\
Наблюдение 15 & 15\\
Наблюдение 16 & 16\\
Наблюдение 17 & 17\\
Наблюдение 18 & 18\\
Наблюдение 19 & 19\\
Наблюдение 20 & 20\\
Наблюдение 21 & 21\\
Наблюдение 22 & 22\\
Наблюдение 23 & 23\\
Наблюдение 24 & 24\\
Наблюдение 25 & 25\\
Наблюдение 26 & 26\\
Наблюдение 27 & 27\\
Наблюдение 28 & 28\\
Наблюдение 29 & 29\\
Наблюдение 30 & 30\\
Наблюдение 31 & 31\\
Наблюдение 32 & 32\\
Наблюдение 33 & 33\\
Наблюдение 34 & 34\\
Наблюдение 35 & 35\\
Наблюдение 36 & 36\\
Наблюдение 37 & 37\\
Наблюдение 38 & 38\\
Наблюдение 39 & 39\\
Наблюдение 40 & 40\\
Наблюдение 41 & 41\\
Наблюдение 42 & 42\\
Наблюдение 43 & 43\\
Наблюдение 44 & 44\\
Наблюдение 45 & 45\\
Наблюдение 46 & 46\\
Наблюдение 47 & 47\\
Наблюдение 48 & 48\\
Наблюдение 49 & 49\\
Наблюдение 50 & 50\\
Наблюдение 51 & 51\\
Наблюдение 52 & 52\\
Наблюдение 53 & 53\\
Наблюдение 54 & 54\\
Наблюдение 55 & 55\\
Наблюдение 56 & 56\\
Наблюдение 57 & 57\\
Наблюдение 58 & 58\\
Наблюдение 59 & 59\\
Наблюдение 60 & 60\\
Наблюдение 61 & 61\\
Наблюдение 62 & 62\\
Наблюдение 63 & 63\\
Наблюдение 64 & 64\\
Наблюдение 65 & 65\\
Наблюдение 66 & 66\\
Наблюдение 67 & 67\\
Наблюдение 68 & 68\\
Наблюдение 69 & 69\\
Наблюдение 70 & 70\\\end{longtable}
\chapter{Проверка следующего раздела}
Ссылка на рисунок~\ref{fig:second} и формулу~\eqref{eq:second}.
\begin{figure}[!htbp]\centering\fbox{Второй раздел}
\caption{Второй рисунок}\label{fig:second}\end{figure}
\begin{equation}\label{eq:second}x=y.\end{equation}
Порядок букв: \GostCyrLetter{1} \GostCyrLetter{2} \GostCyrLetter{3} \GostCyrLetter{4} \GostCyrLetter{5} \GostCyrLetter{6} \GostCyrLetter{7} \GostCyrLetter{8} \GostCyrLetter{9} \GostCyrLetter{10} \GostCyrLetter{11} \GostCyrLetter{12} \GostCyrLetter{13} \GostCyrLetter{14} \GostCyrLetter{15} \GostCyrLetter{16} \GostCyrLetter{17} \GostCyrLetter{18} \GostCyrLetter{19} \GostCyrLetter{20} \GostCyrLetter{21} \GostCyrLetter{22} \GostCyrLetter{23} \GostCyrLetter{24} \GostCyrLetter{25} \GostCyrLetter{26}.
\GostAppendix{Длинный заголовок приложения для проверки переноса строки в содержании и положения продолжения записи относительно обозначения приложения}{справочное}{app:a}
\section{Данные приложения}
Ссылка на рисунок~\ref{fig:app} и формулу~\eqref{eq:app}.
\begin{figure}[!htbp]\centering\fbox{Приложение А}
\caption{Рисунок приложения}\label{fig:app}\end{figure}
\begin{equation}\label{eq:app}u=v.\end{equation}
\GostAppendix{Второе приложение}{справочное}{app:b}
Ссылка на таблицу~\ref{tab:app}.
\begin{table}[!htbp]\caption{Проверка номера}\label{tab:app}
\centering\begin{tabular}{|c|c|}\hline Поле & Значение\\\hline Пример & 1\\\hline\end{tabular}\end{table}
\end{document}
